% Source: https://artofproblemsolving.com/community/c3h3107652_prove_it
% Stable identity: aops-topic-3107652-q01
Suppose a solution exists with least positive $ c $. Dividing out common factors gives $\gcd(a,b)=1 $, and squares modulo $ 16 $ show one of $ a,b $ is even; take $ b $ even. The primitive Pythagorean triple $(a^2,b^2,c)$ has $ a^2=m^2-n^2 $, $ b^2=2mn $, $ c=m^2+n^2 $, with $ m>n $ coprime and of opposite parity. Here $ m $ is odd, for otherwise $ a^2\equiv-1\pmod4 $. Since $ 2mn $ is a square and the factors are coprime, $ m=u^2 $ and $ n=2v^2 $. Therefore $ a^2+(2v^2)^2=u^4 $ is another primitive Pythagorean triple. Parametrizing it gives $ a=r^2-s^2 $, $ v^2=rs $, $ u^2=r^2+s^2 $, where $ r,s $ are coprime positive integers. Consequently $ r=p^2,s=q^2 $, so $ p^4+q^4=u^2 $. But $ 0<u<c=m^2+n^2 $, contradicting the minimal choice. This proves the stronger assertion, and taking $ c^2 $ in place of the hypotenuse gives the fourth-power consequence.
